% !TeX root = main.tex
\subsection*{Incoming design conditions}
The preceding sheet selects C63 protection and a 4-core 35 mm$^2$ feeder with 16 mm$^2$ CPC. Along the meter-to-DB route, the length is 28.5119 m plus a 3 m vertical/termination allowance. Cable sizing assumes 40\,$^\circ$C, a separate Method B enclosure and third-harmonic content greater than 15\% and at most 33\%, giving $C_a=0.87$, $C_g=1$ and $C_h=0.86$. A full-size neutral is retained on this basis.

To bound the phase-neutral voltage drop, the calculation uses $(2/\sqrt3)(1.10)=1.270170592$ mV/A/m from the three-phase table impedance, with neutral current at most 63 A. Thus, feeder drop is 2.521606 V at 63 A and 1.945239 V at the design demand of 48.600 A. Additionally, the complete-path calculation below includes the 2.2 V upstream allowance.

\subsection*{All outgoing cables}
All outgoing cable sizes refer to PVC copper L/N/CPC conductors. Lighting grouped separately has $I_z=24(0.87)(0.54)=11.2752$ A for seven 2.5 mm$^2$ circuits;  all $I_b\le 10\le I_z$. Lighting lengths include 5 m DB/switch vertical allowance.

Each uniform ring bank counts as four current-carrying legs, giving $I_z=41(0.87)(0.65)=23.1855$ A per 6 mm$^2$ leg, above the standard 20 A minimum. With 13 A per point, the ring-network solution gives maximum leg currents 14.658--15.203 A, below $I_z$. However, the radial criterion $I_n\le I_z$ must not be applied to just one ring leg. Continuity and fused-point restrictions are essential.

Meanwhile, the 10 mm$^2$ A2 radial cables have $I_z=57(0.87)(0.70)=34.713$ A;  P7's separate 4 mm$^2$ radial has $I_z=32(0.87)=27.84$ A. In determining route lengths, the power radials include a 2 m vertical allowance, while the rings include 2 m at each DB end. Additionally, each local point spur adds 0.5 m of 4 mm$^2$ cable.
\begin{longtable}{llrrrr}\toprule
Circuit & Phase & $I_b$ A & RCBO A & L/N/CPC mm$^2$ & Main length m\\\midrule\endhead

LC-A & L1 & 2.030 & 10 & 2.5 & 58.301 \\
LC-B & L2 & 1.980 & 10 & 2.5 & 62.184 \\
LC-C & L3 & 2.000 & 10 & 2.5 & 44.999 \\
LC-D & L3 & 1.222 & 10 & 2.5 & 27.233 \\
LC-E & L2 & 1.980 & 10 & 2.5 & 34.464 \\
LC-F & L1 & 1.414 & 10 & 2.5 & 32.710 \\
LC-G & L3 & 0.768 & 10 & 2.5 & 27.379 \\
P1 & L1 & 26.000 & 32 & 6 & 28.620 \\
P2 & L1 & 26.000 & 32 & 6 & 35.520 \\
P3 & L2 & 26.000 & 32 & 6 & 45.220 \\
P4 & L2 & 26.000 & 32 & 6 & 34.220 \\
P5 & L1 & 13.000 & 32 & 10 & 15.607 \\
P6 & L3 & 26.000 & 32 & 10 & 13.864 \\
P7 & L3 & 6.364 & 10 & 4 & 13.107 \\
P8 & L3 & 13.000 & 32 & 10 & 15.107 \\

\bottomrule\end{longtable}
\subsection*{Complete voltage drop and fault checks}
The voltage-drop calculation sums the segment drops using downstream currents on the radials and solves the two-feed ring network, including the stated vertical/local allowances. Combining this result with the 2.2 V pre-origin allowance and the 2.521606 V feeder bound gives the full-path drop, which is checked against a limit of $0.04(220)=8.8$ V.

Fault-protection checks assume a TN-S supply, an external source loop impedance of 0.20 ohm and PFC at most 6 kA at MS-O1/downstream.  Adding the incoming phase+CPC bound of 0.064129 ohm gives a DB external loop impedance of 0.264129 ohm. Table 11(10) reference limits are B10=4.17, B32=1.31 and C63=0.33 ohm. Throughout these checks, the calculation uses operating-temperature paths and a minimum voltage factor of 0.95. Within a uniform ring, the worst parallel path is at most a quarter of the circumference.
\begin{longtable}{lrrrrr}\toprule
Circuit & Out drop V & Total drop \% & $Z_s$ ohm & Min fault A & $Z_s$ limit ohm\\\midrule\endhead

LC-A & 1.450 & 2.805 & 1.314 & 159.1 & 4.17 \\
LC-B & 1.237 & 2.709 & 1.383 & 151.1 & 4.17 \\
LC-C & 0.959 & 2.582 & 1.074 & 194.6 & 4.17 \\
LC-D & 0.422 & 2.338 & 0.754 & 277.1 & 4.17 \\
LC-E & 0.797 & 2.509 & 0.884 & 236.3 & 4.17 \\
LC-F & 0.642 & 2.438 & 0.853 & 245.0 & 4.17 \\
LC-G & 0.345 & 2.303 & 0.757 & 276.1 & 4.17 \\
P1 & 1.269 & 2.723 & 0.322 & 649.3 & 1.31 \\
P2 & 1.675 & 2.907 & 0.334 & 624.9 & 1.31 \\
P3 & 1.937 & 3.027 & 0.352 & 593.5 & 1.31 \\
P4 & 1.428 & 2.795 & 0.332 & 629.4 & 1.31 \\
P5 & 0.964 & 2.584 & 0.338 & 617.8 & 1.31 \\
P6 & 1.310 & 2.742 & 0.331 & 632.1 & 1.31 \\
P7 & 0.953 & 2.579 & 0.414 & 505.1 & 4.17 \\
P8 & 0.936 & 2.571 & 0.336 & 621.8 & 1.31 \\

\bottomrule\end{longtable}
Under the adopted conditions, all 15 full voltage bounds pass, with a worst-case drop of 3.027\%. Additionally, the calculated minimum earth-fault current in each B-curve final circuit exceeds the threshold $5I_n$. Similarly, incoming C63 gives $Z_s=0.264129<0.33$ ohm and $I_{f,\min}=791.28>10(63)=630$ A, satisfying reference instantaneous disconnection under the adopted conditions.
\newpage
\subsection*{Overload and total clearing energy}
CPC thermal withstand is checked using $k=115$ for bunched PVC copper CPCs initially at 70\,$^\circ$C (11(2)(c)), with total $I^2t\le(kS)^2$ required. Additionally, the main protection and RCBOs must have a breaking capacity of 10 kA. B-curve devices must limit class 3-equivalent total energy to at most 70,000 A$^2$s for B10 and 90,000 for B32. Accordingly, the selected RCBOs must have manufacturer total-clearing-energy data that meet these limits. Corresponding CPC withstands are 82,656.25 (2.5 mm$^2$), 211,600 (4), 476,100 (6) and 1,322,500 (10), each above the required device bound.

Incoming C63 total energy must be at most 1,000,000 A$^2$s over the relevant fault range, below CPC16 withstand 3,385,600. Similarly, complete 10/13 A FCU energy must be at most 90,000, below the 4 mm$^2$ spur withstand 211,600. Therefore, the incoming breaker and FCU assemblies must be selected to meet these total-energy limits.

For WH01, $I_b=1400/220=6.364<10<27.84$ A, and the manufacturer specifies 10 A protection. Its 10 A FCU gives $I_2=1.9(10)=19<28.2576$ A for the local spur. The assumed 0.5 m local allowance exceeds the 0.448 m drawn connection. The appliance assembly has an accessible earthed socket and double-pole isolation with at least 3 mm contact separation, following manual section 10.4.

Under socket fuse overload conditions, $I_2=1.9(13)=24.7$ A. Spur $I_z=32(0.87)(0.70)=19.488$ A gives $1.45I_z=28.2576>24.7$ A. Overload suitability of P1--P4 is assessed using the standard A1 arrangement in EMSD 6C, 6E and Table 6(1), Note 1(i): 32 A protection, copper size at least 2.5 mm$^2$, corrected per-leg capacity at least 20 A, served area at most 100 m$^2$, intact L/N/CPC rings and evenly distributed outlets. On this basis, the adopted 6 mm$^2$ rings meet these conditions. Exactly two individually 13 A fused connections are allowed per ring; no additional unfused taps are permitted. Here, the leg-current calculation checks load distribution; the A1 standard arrangement governs ring overload protection. With a reference interruption rating of 6 kA, the fuse meets the assumed PFC of at most 6 kA. However, its typical pre-arcing 870 A$^2$s is \emph{not} total clearing energy. Finally, the design assumes an FCU ambient temperature of at most 35\,$^\circ$C and a board apparatus ambient temperature of at most 30\,$^\circ$C; a containment temperature of 40\,$^\circ$C remains the conservative cable condition.
\subsection*{References}
CLP COP215 Rev 07: clauses 3.1.7 and 3.3.1; Appendices 2 and 7 (building demand). CLP COP101 v16: clauses 4.3, 5.5, 5.6 and 6.1--6.7 (plantroom planning).

EMSD Code of Practice for the Electricity (Wiring) Regulations, 2025: Table 6(1), printed p.56; Table 7(1), pp.65--67; 9B, p.79; Table 11(2), pp.99--100; Table 11(10), p.105; Appendix 5 correction factors and Appendix 6 cable ratings.

Performance in Lighting, PLATO 600 MultiPower 3117955: LOW (22 W), MEDIUMLOW (28 W) and HIGH (38 W) manufacturer LDT files, October 2022. ABB guide 9AKK108468A3912, p.130: B-curve energy references. Eaton TDC180 Technical Data 2042, August 2024, p.2: fuse conventional-fusing and interruption data. Total FCU clearing energy is specified separately above.

\noindent Electrical Schematic 1: transformer, main LV switchboard and floor-metering arrangement. Electrical Schematic 2: final distribution-board schematic conventions.


\clearpage
\subsection*{Selected customer-side apparatus}
\begin{longtable}{p{.19\linewidth}p{.42\linewidth}p{.31\linewidth}}\toprule
Item & Selected model / configuration & Design selection\\\midrule\endhead
Main incomers & Schneider MasterPact MTZ2 25H1, four-pole; MicroLogic 5.0X & 2500 A frame; 2250 A pickup; 66 kA\\
Rising supplies & Schneider MasterPact MTZ2 16H1, four-pole; MicroLogic 5.0X & 1600 A frame; 1440 A pickup; 66 kA\\
Floor protection & Schneider ComPacT NSX400H C40H42D400 / MicroLogic 2.3 & 400 A frame; 320 A pickup; 70 kA\\
Plant/public & Schneider ComPacT NS800H C080H45AFM / MicroLogic 5.0A & 800 A frame; 720 A pickup; 70 kA\\
MS-O1 & Schneider Acti9 iC60H A9F54463 & C63; four-pole; 10 kA Icn\\
LC-A--G and P7 & Schneider Acti9 iC60H A9D31810 & B10; 1P+N; Type A, 30 mA; 10 kA\\
P1--P6 and P8 & Schneider Acti9 iC60H A9D31832 & B32; 1P+N; Type A, 30 mA; 10 kA\\
WH01 via FS01 & STIEBEL ELTRON SHC10 GB Eltron 235232 & 10 L; 1.4 kW at 220 V; B10 and 10 A DP FCU\\
\bottomrule\end{longtable}
The main and rising-supply ACBs require the specified MicroLogic control units; a bare breaker frame is not a complete protective device. The DB retains the specified 100 A four-pole isolator and eight-way TPN arrangement, with 15 installed outgoing devices. Unused positions have blanking plates and no connected loads. The apparatus schedule is an electrical design selection, not evidence of utility approval or installed equipment.

\subsection*{Manufacturer selection references}
\href{https://www.stiebel-eltron.com/content/dam/ste/cdbassets/current/bedienungs-_u_installationsanleitungen/instructionandinstallationmanual_doc-00060782.pdf}{STIEBEL ELTRON, SHC 10 / SHC 15 operation and installation manual}, sections 10.4 and 15.7 (printed p.12): 220 V heater rating, 10 A protection and dimensions. The tabulated rounded current is 6.2 A; this design conservatively uses 1400/220 = 6.364 A.

\href{https://ckm-content.se.com/ckmContent/sfc/servlet.shepherd/document/download/0691H00000FWgC6QAL}{Schneider Electric, MasterPact MTZ circuit breakers overview}: MTZ2 H1 frame range and 66 kA duty at 415 V. Schneider product sheets: \href{https://www.se.com/nl/nl/product/C40H42D400/compact-vermogenschakelaar-nsx400h-70ka-ac-4p4d-400a-2-3/}{C40H42D400} and \href{https://www.se.com/uk/en/product/C080H45AFM/circuit-breaker-compact-ns800h-70ka-at-415vac-4p-fixed-manually-operated-micrologic-5-0a-control-unit-800a/}{C080H45AFM} establish the selected MCCB ratings.

Schneider Acti9 product references \href{https://eshop.se.com/ae/acti9-ic60h-4p-63a-c-miniature-circuit-breaker-a9f54463.html}{A9F54463}, \href{https://www.se.com/uk/en/product/A9D31810/ic60h-earth-leakage-circuit-breaker-1p-\%2B-n-b-curve-10-a-30-ma-240-v/}{A9D31810} and \href{https://www.se.com/uk/en/product/A9D31832/ic60h-earth-leakage-circuit-breaker-1p-\%2B-n-b-curve-32-a-30-ma-240-v/}{A9D31832}. Product coordination must satisfy the declared source and total-energy bounds; full discrimination is not inferred from model numbers.
